% Rewritten from executed/33_diffusion_train.ipynb and pages/33-diffusion_train*.md.

\chapter[Diffusion LM 샘플러 교정]{Diffusion LM 샘플러 교정: carry-over와 반복 억제}
\label{ch:33}
\markboth{Diffusion LM 샘플러 교정}{Diffusion LM 샘플러 교정}

\chaptermeta{33_diffusion_train/33_diffusion_train.ipynb}{https://colab.research.google.com/github/yoon-gu/neuqes-101/blob/master/33_diffusion_train/33_diffusion_train.ipynb}{같은 작은 diffusion LM에서 샘플러만 바꿔 반복을 줄이는 실험}

\index{1carry over sampler@carry-over sampler}
\index{1semi autoregressive@semi-AR}
\index{1carry over semi AR@carry-over semi-AR}
\index{1repetition penalty@repetition penalty}
\index{1repetition suppression@repetition suppression}
\index{1top p@top-p}
\index{1top k@top-k}
\index{1temperature@temperature}
\index{1fixed t accuracy@fixed-t accuracy}
\index{1ngram repetition@4-gram repetition}
\index{1confidence remasking@confidence remasking}
\index{1sampler ablation@sampler ablation}
\index{1unigram collapse@unigram collapse}
\index{1diffusion sampler@diffusion sampler}
\index{1AR vs diffusion@AR vs diffusion}
\index{1parallel denoising@parallel denoising}
\index{1sampling hyperparameter@sampling hyperparameter}
\index{1surrogate loss@surrogate loss}
\index{1one over t weighting@1/t weighting}
\index{1fixed t diagnostic@fixed-t diagnostic}
\index{1four gram repetition rate@4-gram repetition rate}
\index{1carry over state@carry-over state}
\index{1token confirmation@token confirmation}
\index{1mask schedule@mask schedule}
\index{0반복 억제@반복 억제}
\index{0샘플러 교정@샘플러 교정}
\index{0carry over 샘플러@carry-over 샘플러}
\index{0반자기회귀 샘플러@반자기회귀 샘플러}
\index{0반복률@반복률}
\index{0엔그램 반복@엔그램 반복}
\index{0샘플러 ablation@샘플러 ablation}
\index{0토큰 확정@토큰 확정}
\index{0마스크 스케줄@마스크 스케줄}
\index{0고정 t 진단@고정 t 진단}
\index{0대리 손실@대리 손실}
\index{0AR와 diffusion 비교@AR와 diffusion 비교}
\index{0병렬 denoise@병렬 denoise}

\section{누적 추적표}

\begin{booktable}{33장 누적 추적표}{tab:ch33-trace}
\begin{adjustbox}{max width=\textwidth}
\begin{tabular}{@{}lllll@{}}
\toprule
장 & 모델 & 학습 레시피 & 샘플러 & 핵심 결과 \\
\midrule
\ref{ch:32}장 & BERT-MLM 3.79M & BPE 2048, 30,000 step, \(1/t\) loss & confidence-remasking & 작동하지만 반복이 남음 \\
\textbf{33장} & \textbf{동일} & \textbf{동일} & \textbf{carry-over + repetition suppression} & \textbf{4-gram 반복률 0.177 \(\to\) 0.000} \\
\bottomrule
\end{tabular}
\end{adjustbox}
\end{booktable}

\section{변경점: 학습이 아니라 샘플러}

한 챕터 한 변화 원칙에 따라 이 장의 변화는 샘플러 하나입니다. 모델 구조, vocab 2048 BPE, 30,000 step, \(1/t\) 시간가중 loss는 \ref{ch:32}장에서 그대로 이어받습니다. 바뀌는 것은 생성할 때 토큰을 확정하고 반복을 억제하는 방식입니다.

\begin{booktable}{33장 변경점 요약}{tab:ch33-diff}
\begin{adjustbox}{max width=\textwidth}
\begin{tabular}{@{}lll@{}}
\toprule
축 & \ref{ch:32}장 & \ref{ch:33}장 \\
\midrule
모델 & \inlinecode{BertForMaskedLM}, 3.79M & 동일 \\
토크나이저 & BPE vocab 2048 + \inlinecode{[MASK]} & 동일 \\
학습 & 30,000 step, \(1/t\) loss & 동일 \\
샘플러 & confidence-remasking & carry-over semi-AR + 반복 억제 \\
반복 진단 & 반복이 남음 & 4-gram 반복률 0.177에서 0.000 \\
\bottomrule
\end{tabular}
\end{adjustbox}
\end{booktable}

\section{Loss 노트: \(1/t\) 시간가중은 그대로 둡니다}

이 장은 샘플러 장이므로 손실 함수는 \ref{ch:32}장과 같습니다. 흡수형 mask diffusion에서 선형 스케줄 \(\alpha_t=1-t\)를 쓰면 시간가중은 식~\eqref{eq:ch33-timeweight}처럼 정리됩니다.

\begin{equation}
\label{eq:ch33-timeweight}
\left|\frac{\alpha'_t}{1-\alpha_t}\right|
= \frac{1}{t}
\end{equation}

구현에서는 마스크된 자리의 CE 합을 길이 \(L\)로 나눈 뒤 \(1/t\)를 곱합니다. 이 \(/L\) 정규화 때문에 실행 코드가 출력하는 값은 ELBO 자체가 아니라, ELBO와 상수배 관계에 있는 surrogate입니다. 따라서 절대값보다 같은 run 안에서의 학습 추이와 \inlinecode{fixed-t} 진단을 함께 읽어야 합니다.

\section{실습: 같은 모델을 다시 확인합니다}

\begin{lstlisting}
import math, random, time
import torch
from datasets import load_dataset
from tokenizers import ByteLevelBPETokenizer
from transformers import BertConfig, BertForMaskedLM

device = "cuda" if torch.cuda.is_available() else "cpu"
USE_FP16 = torch.cuda.is_available()
print(f"torch {torch.__version__} | device {device} | fp16 {USE_FP16}")
if torch.cuda.is_available():
    print("GPU:", torch.cuda.get_device_name(0))
\end{lstlisting}

\begin{lstlisting}[style=bookoutput]
torch 2.11.0+cu128 | device cuda | fp16 True
GPU: Tesla T4
\end{lstlisting}

\begin{lstlisting}[style=bookoutput]
vocab_size : 2048
mask_id    : 2 | pad_id: 0
sample tok : ['<sp>Once', '<sp>upon', '<sp>a', '<sp>time',
              '<sp>there', '<sp>was', '<sp>a', '<sp>little',
              '<sp>cat', '.']
\end{lstlisting}

\noindent\emph{출력 해석.}
실행본의 ByteLevel BPE 토큰은 공백 표시 기호를 포함합니다. 원고에서는 폰트 누락을 피하기 위해 그 기호를 \inlinecode{<sp>}로 표기했습니다. 의미는 “앞에 공백이 붙은 토큰”입니다.

\begin{lstlisting}[style=bookoutput]
#params 3.79M | embedding share 13.9%

=== summary ===
elapsed 18.50 min | step 30000 | train_loss 3.5916
random baseline ln(V) = 7.6246
peak VRAM 627 MiB
\end{lstlisting}

\textbf{출력 해석.}
최종 \inlinecode{train\_loss}는 3.5916입니다. \ref{ch:32}장의 3.7148과 같은 규모이며, 본 장의 품질 차이는 loss가 아니라 샘플러에서 확인해야 합니다.

\section{샘플러 설계}

기본 confidence-remasking은 병렬 복원의 장점을 살리지만, 이미 확정된 반복을 잘 깨지 못합니다. 본 장의 샘플러는 세 가지를 추가합니다.

\begin{itemize}
\tightlist
\item carry-over: 이전 step에서 확정한 토큰을 다음 step 입력으로 넘깁니다.
\item nucleus sampling: \inlinecode{temperature=0.8}, \inlinecode{top\_p=0.92}로 너무 넓은 후보를 줄입니다.
\item 반복 억제: \inlinecode{repetition\_penalty=1.3}과 인접 동일 토큰 금지로 짧은 반복을 직접 막습니다.
\end{itemize}

\begin{lstlisting}
def suppress_repetition(logits, generated, penalty=1.3):
    for token_id in set(generated):
        logits[token_id] /= penalty
    if generated:
        logits[generated[-1]] = -float("inf")
    return logits

def sample_with_top_p(logits, temperature=0.8, top_p=0.92):
    logits = logits / temperature
    probs = torch.softmax(logits, dim=-1)
    sorted_probs, sorted_idx = torch.sort(probs, descending=True)
    keep = torch.cumsum(sorted_probs, dim=-1) <= top_p
    keep[..., 0] = True
    filtered = torch.zeros_like(probs)
    filtered.scatter_(0, sorted_idx[keep], sorted_probs[keep])
    filtered = filtered / filtered.sum()
    return torch.multinomial(filtered, 1).item()
\end{lstlisting}

\noindent\emph{위 코드 읽기.}
\inlinecode{suppress\_repetition}은 이미 나온 토큰의 logit을 낮추고, 직전 토큰은 다시 뽑히지 못하게 막습니다. \inlinecode{sample\_with\_top\_p}는 확률 질량의 상위 92\% 후보에서만 샘플링합니다. 이 조합은 모델을 바꾸지 않고 decoding만 바꾸는 실험입니다.

\section{정량 진단}

샘플러 변경 효과는 두 지표로 봅니다. 고정-\(t\) 복원 정확도는 모델이 중간 노이즈 상태에서 원래 토큰을 얼마나 맞히는지 보고, 4-gram 반복률은 생성 결과의 짧은 반복을 봅니다.

\begin{lstlisting}[style=bookoutput]
[diag] fixed-t(0.15) top-1 accuracy = 0.717
[diag] 4-gram repeat rate  A(default)=0.177  B(rep-suppressed)=0.000
\end{lstlisting}

\textbf{출력 해석.}
고정-\(t=0.15\)에서 top-1 accuracy는 0.717입니다. 모델 자체가 복원 문제를 어느 정도 풀고 있음을 보여줍니다. 같은 모델에서 4-gram 반복률이 0.177에서 0.000으로 내려갔으므로, 본 장의 개선은 학습보다 샘플러의 반복 억제에서 온다고 해석할 수 있습니다.

\section{생성 비교}

아래 출력은 같은 학습 모델과 같은 프롬프트에서 샘플러만 바꾼 결과입니다.

\begin{lstlisting}[style=bookoutput]
=== sampler sweep (same model, prompt 'Once upon a time') ===

----- A) default temp0.7/topk40 -----
[0]  Once upon a time, there was a little girl named Lily. She loved
to play with her toy friends. One day, Lily's mom came to play with her
toys. She saw a big ball and wanted to play with it.

Lily's mom asked her to help her mom. She asked her mom if she could
play with her toy ball. Lily said, "Yes, you can play with my ball with it."

----- B) rep1.3 + no adjacent repeat + topp0.92 -----
[0]  Once upon a time, there was a little girl named Lily. She loved
to play with her friends and explore in her backyard. One day, she went
outside to the park when she saw a big ball! It looked so pretty that it
wanted to catch it.

Lily picked up and tried to grab it inside. But then, she found a rock
on the ground and started to run away.
\end{lstlisting}

\textbf{출력 해석.}
A는 \inlinecode{ball with it}처럼 같은 의미 단위가 가까운 거리에서 반복됩니다. B는 완벽한 이야기는 아니지만 같은 짧은 구를 되풀이하는 현상이 줄어듭니다. 이 차이가 4-gram 반복률 0.177 \(\to\) 0.000과 맞물립니다.

\section{오류 실험: 레시피를 되돌리면 왜 붕괴하나}

본 장의 중심은 샘플러입니다. 다만 교육적으로, 레시피를 일부러 되돌렸을 때 생기는 붕괴도 보존합니다. vocab을 너무 크게 잡거나, step을 줄이거나, \inlinecode{[MASK]} 처리를 느슨하게 하면 모델은 random baseline 근처에 머물거나 \inlinecode{[MASK]}를 그대로 출력하기 쉽습니다.

\begin{booktable}{33장 붕괴 ablation 해석}{tab:ch33-ablation}
\begin{adjustbox}{max width=\textwidth}
\begin{tabular}{@{}lll@{}}
\toprule
되돌린 요소 & 관찰되는 현상 & 해석 \\
\midrule
큰 vocab & embedding 비중이 커짐 & 작은 모델에서 단어장 파라미터가 본체를 압도 \\
짧은 학습 & loss가 random baseline 근처 & diffusion 복원은 충분한 step이 필요 \\
단순 샘플러 & 짧은 구 반복 & 복원 정확도와 생성 품질 사이에 decoding gap 존재 \\
\bottomrule
\end{tabular}
\end{adjustbox}
\end{booktable}

\section{AR 모델과 diffusion 모델 비교}

\begin{booktable}{33장 AR과 diffusion 생성 비교}{tab:ch33-ar-diffusion}
\begin{adjustbox}{max width=\textwidth}
\begin{tabular}{@{}lll@{}}
\toprule
항목 & Autoregressive LM & Diffusion LM \\
\midrule
생성 단위 & 다음 토큰 하나 & 여러 위치 병렬 복원 \\
학습 신호 & \(\log p(x_i \mid x_{<i})\) & \(\log p(x_i \mid \tilde{x}_t)\) \\
장점 & 간단하고 강력한 기준선 & step 수와 순서를 설계할 여지 \\
약점 & 순차 생성 비용 & 샘플러·스케줄 설계가 품질을 좌우 \\
\bottomrule
\end{tabular}
\end{adjustbox}
\end{booktable}

\begin{faqBox}{Q1. 본 장은 \ref{ch:32}장의 붕괴를 고친 장인가요?}
아닙니다. 최신 \ref{ch:32}장은 이미 작동합니다. 본 장은 같은 모델에서 샘플러만 바꿔 반복을 줄이는 장입니다. 붕괴 내용은 레시피를 일부러 되돌린 ablation으로만 다룹니다.
\end{faqBox}

\begin{faqBox}{Q2. \inlinecode{train\_loss} 3.5916과 반복률 0.000 중 무엇이 더 중요한가요?}
목적이 생성 품질이면 둘 다 봐야 합니다. loss는 복원 학습이 됐는지 확인하는 지표이고, 반복률은 decoding 결과의 품질을 보는 지표입니다.
\end{faqBox}

\begin{faqBox}{Q3. 반복 억제는 모델 성능을 속이는 것 아닌가요?}
샘플러는 생성 모델의 일부입니다. 같은 language model이라도 temperature, top-p, repetition penalty에 따라 실제 출력 품질이 크게 달라집니다. 다만 샘플러가 모델의 지식 자체를 늘리지는 않습니다.
\end{faqBox}

\begin{faqBox}{Q4. carry-over는 autoregressive 생성과 같나요?}
완전히 같지는 않습니다. AR은 왼쪽에서 오른쪽으로 한 토큰씩 조건을 누적합니다. carry-over diffusion은 여러 위치를 병렬로 복원하되, 이전 step에서 확정한 토큰을 다음 step 조건으로 유지합니다.
\end{faqBox}

\begin{previewBox}{미리보기: 다음 장}
\ref{ch:34}장에서는 같은 diffusion 아이디어를 한국어로 옮깁니다. 단순히 데이터만 바꾸면 unigram collapse가 나타나고, BERT의 80/10/10 마스킹이 왜 안정화에 도움이 되는지 확인합니다.
\end{previewBox}
